\chapter{Dicas Gerais Relacionadas ao LaTeX}

\section{Citações}

Para inserir citações em um documento LaTeX, é necessário primeiro salvar as informações de cada fonte em um arquivo de extensão .bib, conhecido como arquivo BibTeX. Este arquivo contém as referências bibliográficas no formato BibTeX, que é utilizado pelo LaTeX para gerar automaticamente as citações e a lista de referências no documento final. 

Neste modelo, há um arquivo \verb|bibliografia.bib|  que deve ser mantido no mesmo diretório que o arquivo principal do LaTeX (\verb|main.tex|), e cada entrada deve seguir a sintaxe apropriada para garantir a correta formatação das~referências.

Abaixo são mostrados alguns exemplos de referências em formato BibTeX:

\begin{lstlisting}[caption={Exemplos de entradas BibTeX},label={lst:codigo-bibtex},basicstyle=\footnotesize\ttfamily]
% LIVRO
@book{gil2008metodos,
  title={Métodos e Técnicas de Pesquisa Social},
  author={Gil, Antonio Carlos},
  year={2008},
  address={São Paulo},
  publisher={Atlas}
}

% ARTIGO
@article{Gautschi2010espiral,
    author = {Walter Gautschi},
    title = {The spiral of Theodorus, numerical analysis, and special functions},
    journal = {Journal of Computational and Applied Mathematics},
    volume = {235},
    number = {4},
    pages = {1042-1052},
    year = {2010},
}

% SITE
@Misc{wolfram2024espiral,
    author = {Symon Tyler},
    title = {The Spiral of Theodorus},
    year = {2024},
    url = {https://demonstrations.wolfram.com/TheSpiralOfTheodorus/},
    urlaccessdate = {03 ago. 2024}
}
\end{lstlisting}

A inclusão do link deve ser opcional. Obras disponíveis apenas em nossas bases, sem acesso público, não devem receber link.

Exemplo de citação com página \cite[p. 1]{gil2008metodos} e sem página \cite{gil2008metodos}.

Exemplo exibindo somente o ano: \citeyear{gil2008metodos}.

Exemplo exibindo somente o autor: \citeauthor{gil2008metodos}.

Exemplo de citação integrada ao texto com página: \citeonline[p. x]{gil2008metodos} realizou...


\section{Listas}

O LaTeX permite criar listas numeradas utilizando o ambiente \texttt{enumerate}.

\subsection{Exemplo}
  
Neste modelo, configuramos o ambiente para exibir \textbf{a), b), c)} e também subalínea usando o ambiente \texttt{itemize}:

\begin{enumerate}
  \item Primeiro objetivo;
      \begin{itemize}
          \item Subalínea A;
          \item Subalínea B.
      \end{itemize}
  \item Segundo objetivo;
  \item Terceiro objetivo.
\end{enumerate}

\section{Notas de Rodapé}

Em alguns casos, pode ser necessário adicionar informações complementares sem interromper o fluxo principal do texto. Para isso, o LaTeX oferece o comando \verb|\footnote{}| para notas de rodapé.

\subsection{Exemplo}

Exemplo de nota de rodapé\footnote{Texto da nota de rodapé de exemplo.}. O resultado\footnote{Outra nota de rodapé.} será que, ao final da página, aparecerá o número da nota seguido do texto correspondente.

\subsection{Boas práticas}
\begin{enumerate}
  \item Use notas de rodapé apenas para informações secundárias ou referências adicionais que não merecem um parágrafo inteiro;
  \item Evite excesso de notas de rodapé, pois podem prejudicar a leitura;
  \item O comando \verb|\footnote{}| pode ser usado em qualquer lugar dentro do corpo do texto, inclusive dentro de ambientes como \verb|itemize| ou \verb|enumerate|.
\end{enumerate}



\section{Figuras}

Figuras são elementos visuais importantes em trabalhos acadêmicos, usados para ilustrar conceitos, resultados ou exemplos práticos.  
No LaTeX, figuras são inseridas dentro do ambiente \verb|figure| (ou, neste template, com o ambiente adaptado \verb|figura|).

\subsection{Exemplo}

O código abaixo insere uma figura com legenda e indicação de fonte, usando o ambiente \verb|figura|:

\begin{verbatim}
\begin{figura}{0.8\textwidth}
    \includegraphics[width=\textwidth]{img/exemplo.png}
    \caption{Exemplo de figura inserida em \LaTeX.}
    \label{fig:exemplo}
    \sourceFig{os autores, 2025.}
\end{figura}
\end{verbatim}

\noindent Resultado:

\begin{figura}{0.6\textwidth}
    \caption{Exemplo de figura inserida em LaTeX.}
    \includegraphics[width=\textwidth]{example-image}
    \label{fig:exemplo}
    \sourceFig{os autores, 2025.}
\end{figura}

\subsection{Boas práticas}
\begin{enumerate}
  \item Sempre adicione uma legenda clara e objetiva usando \verb|\caption{}|;
  \item Inclua a fonte da figura (mesmo quando for de autoria própria) usando o comando \verb|\sourceFig{}|;
  \item Nomeie a figura com \verb|\label{}| para poder referenciá-la no texto com \verb|\ref{}|;
  \item Centralize a figura para melhor apresentação, mas ajuste sua largura (ex.: \verb|0.8\textwidth|) para que não ultrapasse as margens do documento.
\end{enumerate}

\section{Gráficos}

A fim de permitir a inclusão de gráficos, bem como que estes sejam apresentados em uma ``Lista de Gráficos'', o grupo de trabalho AbnTeX2 desenvolveu um ambiente customizado, denominado \verb|grafico|, baseado no ambiente de base LaTeX~\verb|figure|. 

\subsection{Exemplo}

O código abaixo insere um gráfico com legenda e indicação de fonte, usando o ambiente \verb|grafico|:

\begin{verbatim}
\begin{grafico}{0.8\textwidth}
    \includegraphics[width=\textwidth]{img/exemplo.png}
    \caption{Exemplo de figura inserida em \LaTeX.}
    \label{gra:exemplo}
    \sourceFig{as autoras, 2025.}
\end{grafico}
\end{verbatim}

\noindent Resultado:

\begin{grafico}{0.6\textwidth}
    \caption{Exemplo de gráfico inserido em LaTeX.}
    \includegraphics[width=\textwidth]{example-image}
    \label{gra:exemplo}
    \sourceFig{as autoras, 2025.}
\end{grafico}

\section{Tabelas}

As tabelas permitem organizar dados de forma estruturada e comparativa.  
No LaTeX, tabelas são inseridas dentro do ambiente \verb|table| (ou, neste template, com o ambiente adaptado \verb|tabela|).

\subsection{Exemplo com \texttt{\textbackslash hline}}

O código abaixo insere uma tabela com legenda e indicação de fonte, usando o ambiente \verb|tabela| e linhas horizontais simples:

\begin{verbatim}
\begin{tabela}{0.45\textwidth}
    \caption{Exemplo de tabela inserida em LaTeX.}
    \label{tab:exemplo}
    \renewcommand{\arraystretch}{1.1} % aumenta a altura das linhas
    \resizebox{\textwidth}{!}{
    \begin{tabular}{lcc}
        \hline
        Curso        & Alunos & Percentual \\
        \hline
        Administração & 24    & 32,0\% \\
        Engenharia    & 31    & 41,3\% \\
        Letras        & 10    & 13,3\% \\
        Pedagogia     & 10    & 13,3\% \\
        \hline
    \end{tabular}
    }
    \sourceTab{o autor, 2025.}
\end{tabela}
\end{verbatim}

\noindent Resultado:

\begin{tabela}{0.45\textwidth}
  \caption{Exemplo de tabela inserida em LaTeX.}
  \label{tab:exemplo}
  \renewcommand{\arraystretch}{1.1} % aumenta a altura das linhas
  \resizebox{\textwidth}{!}{
  \begin{tabular}{lcc}
    \hline
    Curso        & Alunos & Percentual \\
    \hline
    Administração & 24    & 32,0\% \\
    Engenharia    & 31    & 41,3\% \\
    Letras        & 10    & 13,3\% \\
    Pedagogia     & 10    & 13,3\% \\
    \hline
  \end{tabular}
  }
  \sourceTab{o autor, 2025.}
\end{tabela}

\subsection{Exemplo com o pacote \texttt{booktabs}}

Também é possível melhorar a estética da tabela utilizando o pacote \verb|booktabs|, que fornece os comandos \verb|\toprule|, \verb|\midrule| e \verb|\bottomrule| para substituir o uso de \verb|\hline|:

\begin{verbatim}
\begin{tabela}{0.45\textwidth}
    \centering
    \caption{Exemplo (2) de tabela inserida em LaTeX.}
    \label{tab:exemplo-booktabs}
    \resizebox{\textwidth}{!}{
    \begin{tabular}{lcc}
        \toprule
        Curso         & Alunos & Percentual \\
        \midrule
        Administração & 24     & 32,0\% \\
        Engenharia    & 31     & 41,3\% \\
        Letras        & 10     & 13,3\% \\
        Pedagogia     & 10     & 13,3\% \\
        \bottomrule
    \end{tabular}
    }
    \sourceTab{a autora, 2025.}
\end{tabela}
\end{verbatim}

\noindent Resultado:

\begin{tabela}{0.45\textwidth}
    \centering
  \caption{Exemplo (2) de tabela inserida em LaTeX.}
  \label{tab:exemplo-booktabs}
  \resizebox{\textwidth}{!}{
  \begin{tabular}{lcc}
    \toprule
    Curso         & Alunos & Percentual \\
    \midrule
    Administração & 24     & 32,0\% \\
    Engenharia    & 31     & 41,3\% \\
    Letras        & 10     & 13,3\% \\
    Pedagogia     & 10     & 13,3\% \\
    \bottomrule
  \end{tabular}
  }
  \sourceTab{a autora, 2025.}
\end{tabela}

\subsection{Boas práticas}
\label{tab:boas_praticas}
\begin{enumerate}
  \item Mantenha as legendas sempre acima das tabelas, de acordo com o padrão acadêmico mais comum;
  \item Utilize o comando \verb|\sourceTab{}| para indicar a fonte dos dados (mesmo que seja o próprio autor);
  \item Nomeie as tabelas com \verb|\label{}| para poder referenciá-las no texto com \verb|\ref{}|;
  \item Evite excesso de linhas horizontais ou verticais; prefira tabelas limpas — o pacote \verb|booktabs| ajuda nisso;
  \item Limite o número de colunas e ajuste a largura da tabela (ex.: \verb|0.45\textwidth|) para garantir boa legibilidade.
\end{enumerate}

\section{Quadros}

A fim de permitir a inclusão de quadros, bem como que estes sejam apresentados em uma ``Lista de Quadros'', o grupo de trabalho AbnTeX2 desenvolveu um ambiente customizado, denominado \verb|quadro2|, baseado no ambiente de base LaTeX~\verb|table|. 

Veja um exemplo no Quadro \ref{qua:exemplo-quadro}.


\begin{quadro2}{0.8\textwidth}
    \caption{Exemplo de quadro inserido em LaTeX.}
    \label{qua:exemplo-quadro}
    \renewcommand{\arraystretch}{1.2} % aumenta a altura das linhas
    \resizebox{\textwidth}{!}{
        \begin{tabular}{|p{6cm}|p{10cm}|}
            \hline
            \textbf{Revista/Publicação} & \textbf{Descrição}\\ \hline
            Revista do Professor de Matemática (RPM) & A RPM é uma publicação voltada para educadores de Matemática, especialmente nas etapas finais do ensino fundamental e no ensino médio. A revista divulga artigos de níveis elementares a avançados, acessíveis tanto para professores de ensino médio quanto para alunos de Licenciatura em Matemática.  \\ \hline
            %
            Matemática Universitária & A Matemática Universitária, publicação semestral da Sociedade Brasileira de Matemática, visa divulgar ideias e estimular a curiosidade intelectual no ensino e estudo da Matemática em nível superior. Destina-se a professores, pesquisadores e alunos de graduação e pós-graduação, promovendo o intercâmbio entre os membros da comunidade acadêmica. \\ \hline
            %
            Revista Professor de Matemática Online (PMO) & A revista Professor de Matemática Online (PMO) publica e divulga artigos acadêmicos relevantes para a formação inicial e continuada de professores da Educação Básica, abordando temas de Matemática, práticas de ensino, história e aplicações. Aceita também resultados de trabalhos de conclusão de curso, ferramentas virtuais e outros produtos de docentes e discentes dos cursos de formação de professores de Matemática. Os artigos são apresentados em português, com resumos em português e inglês. \\ \hline
            %
            Ensino de Matemática em Debate & Discussão sobre metodologias e práticas no ensino de matemática.\\ \hline
        \end{tabular}
    }
    \sourceTab{o autor, 2025.}
\end{quadro2}

\section{Referências Cruzadas}

Você pode criar referências cruzadas para seções, figuras, gráficos, tabelas, quadros e outros elementos usando os comandos \texttt{\textbackslash label} e \texttt{\textbackslash ref}.

No exemplo de inserção de quadros apresentado anteriormente, foi utilizado o comando \verb|\label{qua:exemplo-quadro}|, em que \verb|qua:exemplo-quadro| funciona como um identificador único do elemento. Para referenciá-lo no texto, basta escrever algo como \verb|Conforme Quadro~\ref{qua:exemplo-quadro}|, que produzirá o resultado desejado: Conforme Quadro~\ref{qua:exemplo-quadro}.